% Introduction Document
\section{Introduction}Computational design and analysis has become a fundamental part of industry and academia for use in research, development, and manufacturing. In general, the accuracy of a computational analysis depends heavily on the fidelity of the computational representation of a real-world object or phenomenon. However, the task of creating high fidelity models of an actual geometry can be time-consuming--sometimes consuming up to seventy-five percent of the time required to produce a solution \cite{bischoff05}. This work seeks to improve the process of creating a computational model of an object of interest by accelerating the process of mesh generation. Below in Figure-\ref{GridGenerationProcess} an example of a mesh generation hierarchy can be seen. In general, a valid volume grid (three-dimensional) is bounded by surface grids (two-dimensional); and surface grids are bounded by edge-grids (one-dimensional). At the start of the grid generation hierarchy are point spacing values at the end points of analytical or parametric curves—which bound the edge-grids. In most cases, volume-grid generation is a highly automated procedure—once the bounding surfaces are generated. The same generalization can be made for surface-grids and edge-grids. Algorithms that use automated point creation/insertion for mesh generation of one-, two-, and three-dimensional geometries are ubiquitous [REFERENCES, CUBIT, Delaunay, AFLR]. However the authors are only aware of one such automated scheme for setting point spacing values for more automated mesh generation \cite{mclaurin12}.\\

\begin{figure}[h!]
 \center{\includegraphics[width=0.8\textwidth]
  {Figures/GridGenProcess.eps}}
 \caption{\label{GridGenerationProcess} Specified point spacings and final point distribution for a surface edge \cite{thompson98}}
\end{figure}

The computation involved with edge-grid generation is trivial when compared to volume-grid generation—even with high-order NURBS curves. However, the point spacing values at the end points of curves have to be set manually in order to satisfy a desired length scale. This manual process is time consuming. If geometry repair (gluing, trimming, etc…) is not considered, the amount of user input required to generate a volume grid can be concentrated on the lowest levels of the grid generation hierarchy. In addition, if the edge-grids are not generated appropriately then the errors present there, such as overly dense or sparse spacing, will be propagated up the hierarchy and be present in each subsequent higher-dimensional entity. In order to accelerate the process of generating edge grids, some degree of automation is needed to help deal with common issues associated with undesirable edge-grids.

The proposed algorithm is a general-use method that can be applied to any ``digital curve'' regardless of its representation. This is due to every step being developed without the use of derivatives. Most other methods operate on a specific type of curve, such as NURBS or B-splines, and use the specific information available for the type of curve in use. While in the field of CAD NURBS curves are the de-facto standard, in other fields, such as pattern-recognition, other types of digital curves, such as parametric, are more common[REFERENCES]. T-splines are also becoming more popular in applications such as isogeometric analysis [REFERENCES, Thomas Hughes]. Therefore a general algorithm that does not require derivative information to generate a suitable edge grid was developed. A result of not using derivatives is that each step in the algorithm is robust to large changes in derivatives or curves that are not ``well behaved''---highly oscillatory for instance.

The justification for the development of these methods lies in the need for an automated way of setting point spacing values on curves. This process can only be automated if some way of judging ``how well'' an edge grid represents a curve is present. To this end, two related methods of generating edge grids through constrained optimization are detailed below. Further discussion of element quality, robustness, and a framework for implementing the information associated with an optimal edge grid into an existing grid generator is also presented. Generating edge grids in a more automated fashion accelerates the process of surface grid generation---and ultimately volume grid generation. Using the detailed methods, or any other automated method for setting point spacings does not change the number of steps required for grid generation. However, it does reduce the number of manual steps involved in starting the process.

